\documentclass[11pt,a4paper]{article}
\usepackage[margin=1in]{geometry}
\usepackage{amsmath,amssymb,amsthm}
\usepackage{booktabs}
\usepackage[hidelinks,breaklinks]{hyperref}
\usepackage{enumitem}
\theoremstyle{plain}\newtheorem{theorem}{Theorem}\newtheorem{proposition}[theorem]{Proposition}
\newtheorem{lemma}[theorem]{Lemma}\newtheorem{corollary}[theorem]{Corollary}
\theoremstyle{definition}\newtheorem{example}[theorem]{Example}
\theoremstyle{remark}\newtheorem{remark}[theorem]{Remark}
\newcommand{\N}{\mathbb{N}}\newcommand{\Z}{\mathbb{Z}}
\title{Realizer Bounds and Repetition Constraints\\ at the First Coefficient Crossing in the $3x+1$ Problem}
\author{Elias De Jes\'us\\ \textit{Independent Researcher}\\
ORCID \href{https://orcid.org/0009-0007-0190-9143}{0009-0007-0190-9143}}
\date{14 November 2026}
\begin{document}\maketitle

\begin{abstract}
Terras's coefficient stopping time conjecture asks whether the shortcut Collatz orbit of every
integer $n\ge2$ lies below $n$ when its coefficient $\rho_k=3^d/2^k$ first falls below one.
Here $k$ counts shortcut steps and $d$ counts their odd inputs. We prove that every eligible
realizer descends when the first-crossing word has at most $75$ maximal runs of odd inputs.
Any counterexample requires at least $76$ such runs, $10\,439\,860\,591$ steps and
$6\,586\,818\,670$ odd inputs. The proof uses a lower bound on the seed from the length of its
odd runs, rational approximations to $\log_2 3$, the verified range of Rozier and Terracol,
and Rhin's estimate for a linear form in logarithms. We also prove descent for a family with
arbitrarily many odd runs. Two further restrictions narrow the remaining words. An initial
square $UU$ in a failing word satisfies $|U|<35\log_2 k$. For a counterexample of shortest
crossing length, the reciprocal-sum estimate of Nikzad, Frinkle and Bagherzadeh gives a lower
bound on $\rho_k$ whose value begins $0.9441733669\ldots$. The open inequality compares
the least eligible starting value in each remaining parity class with an explicit threshold.
\end{abstract}

\section{What is being asked}
Write $T(n)=(3n+1)/2$ at odd inputs and $T(n)=n/2$ at even inputs. Fix $n\ge2$.
For the parity word $w$ of its first $k$ input states $n,T(n),\dots,T^{k-1}(n)$, let $d$
count the odd inputs. Induction gives
\begin{equation}\label{eq:affine}
2^kT^k(n)=3^dn+C_w,
\end{equation}
where the nonnegative integer $C_w$ depends only on $w$. Write $\rho_k=3^d/2^k$ with
$\rho_0=1$. An odd step multiplies the coefficient by $3/2$ and an even step halves it.
Let $\tau(n)$ be the first $k$ with $\rho_k<1$, and let $\chi(n)$ be the first $k$
with $T^k(n)<n$. An empty defining set gives stopping time $\infty$.
Equation \eqref{eq:affine} gives $\tau(n)\le\chi(n)$, so any failure of Terras's conjecture
$\tau(n)=\chi(n)$ has finite $k=\tau(n)$ and $T^k(n)\ge n$
\cite{terras1976,hikawa2026}. An even starting value descends at its first step, so a failure
has odd $n$ and $d\ge1$.

A first-crossing word has $\rho_j\ge1$ for every $j<k$ and $\rho_k<1$. Its last step is
even, so $\rho_k\in[1/2,1)$. With $d\ge1$, this also gives
$2^{k-1}\le 3^d<2^k$ and $k=\lceil d\log_2 3\rceil\le2d$.
An \emph{odd run} is a maximal block of consecutive odd input states.

\section{The signed comparator}\label{sec:comparator}
Dividing \eqref{eq:affine} by $2^k$ gives the form we use throughout.
\begin{proposition}\label{prop:comparator}
With $M_w=C_w/3^d$,
\[
T^k(n)=\rho_k\,(n+M_w),\qquad\text{so}\qquad T^k(n)<n\iff n\bigl(\rho_k^{-1}-1\bigr)>M_w .
\]
Moreover, if the odd runs have lengths $\ell_1,\dots,\ell_s$ and run $i$ starts at time $t_i$,
\begin{equation}\label{eq:runsum}
M_w=\sum_{i=1}^{s}\frac{1-(2/3)^{\ell_i}}{\rho_{t_i}} .
\end{equation}
\end{proposition}
\begin{proof}
The first display is \eqref{eq:affine}. For \eqref{eq:runsum}, each odd step at time $j$
contributes $2^j3^{-d_{j+1}}=1/(3\rho_j)$ to $M_w$, where $d_{j+1}$ counts odd steps through time
$j$. Summing a run of $\ell$ consecutive odd steps from $t$ gives
$\frac1{3\rho_t}\sum_{r<\ell}(2/3)^r=(1-(2/3)^\ell)/\rho_t$.
\end{proof}
If a first-crossing word contains $s\ge1$ odd runs, every proper prefix has $\rho_t\ge1$,
so each run contributes less than one and each odd step contributes at most $1/3$. Thus
\begin{equation}\label{eq:Mbounds}
M_w<s,\qquad M_w\le d/3 .
\end{equation}
The two bounds depend respectively on the number of runs and on the total number of odd inputs.

\begin{remark}[Testing all eligible realizers of a word]
A length-$k$ parity word specifies one residue class modulo $2^k$.
Indeed, the two lifts of a given length-$j$ class modulo $2^{j+1}$ have starting values that
differ by $2^j$. Their images after $j$ steps differ by $3^{d_j}$, an odd integer, so exactly
one lift has the prescribed next parity. For a first-crossing word, let $r(w)$ be the least
integer $\ge2$ in its class. Every eligible realizer descends if and only if
$r(w)(\rho_k^{-1}-1)>M_w$, since the left side increases with the seed.
If a seed $n$ fails, its endpoint satisfies $0\le T^k(n)-n<s$ by
Proposition \ref{prop:comparator} and \eqref{eq:Mbounds}.
\end{remark}

\begin{remark}[What is classical here]\label{rem:classical}
Almost all of the above is in Lagarias's survey \cite[\S2.2--\S2.3]{lagarias1985}, and the
reader should not take any of it as new.

The residue-class correspondence is Terras \cite[Thm.~1.2, Cor.~1.3]{terras1976} and is
Theorem~B of \cite{lagarias1985}; Garner \cite[\S4]{garner1981} gives it in the
coefficient-stopping-time form, modulo $2^{1+p(m)}$. Lagarias's equations (2.4)--(2.6) are
Proposition~\ref{prop:comparator}: his $T^{(k)}(n)=\lambda_k(n)n+\rho_k(n)$ has
$\lambda_k=\rho_k$ in our notation and $\rho_k(v)=C_w/2^k$, so that $M_w=\rho_k(v)/\lambda_k$.
His (2.13)--(2.14) then state that the elements of a class with
\[
n\;>\;\varepsilon_k^{-1}\rho_k(v),\qquad \varepsilon_k=1-\lambda_k(v),
\]
have $\omega(n)=\sigma(n)$, and rearranging that inequality gives
$n(\rho_k^{-1}-1)>M_w$ \emph{exactly}. \textbf{The descent criterion above is therefore
Lagarias's, not ours.}

Sharper still, Theorem~E of \cite{lagarias1985} states that there is an effectively computable
$k_0$ such that for admissible $\mathbf v$ of length $k\ge k_0$, every element of $S(\mathbf
v)$ except possibly the smallest has stopping time $k$. Its proof is the one we use: a linear
forms in logarithms bound (there Baker--Feldman, $|k\log2-l\log3|\ge k^{-c_0}$; here Rhin's
$k^{-13.3}$), fed into (2.14). \textbf{``Only the least element of the class can fail'' is thus
Theorem~E, asymptotically in $k$.}

What the present section adds is only this. The criterion is stated with no ``sufficiently
large'' and no ineffective $k_0$: it is a single monotone comparison valid at every $k$, so one
evaluation at $r(w)$ decides the class outright. And a failure is pinned to the $s$ endpoints
$n,\dots,n+s-1$. The quantitative work is in \S\ref{sec:75}, where the run decomposition
\eqref{eq:runsum} replaces Lagarias's cruder admissibility estimate
$\rho_k(\mathbf v)\le k2^{(1-\theta)k}$ and makes $k_0$ explicit.
\end{remark}

\begin{example}\label{ex:59}
The exact odd-to-odd valuation word $(1,2,1,3,1)$ from
\cite[App.~A.6(b)]{dj-floors2026} has eight shortcut input bits $11011001$ and an odd
terminal state. The eight input bits give the class $187\bmod 256$. Terminal oddness adds a
ninth bit and selects $187\bmod 512$. Indeed, $T^8(187)=179$, while the other lift, $443$,
has $T^8(443)=179+3^5=422$. The first coefficient crossing already occurs at step $7$.
Its prefix $1101100$ gives the class $59\bmod128$, and
\[
 59-T^7(59)=\frac{(2^7-3^4)\cdot59-85}{2^7}=21.
\]
Thus $T^7(59)=38$ and every eligible seed in $59+128\Z$ descends at step $7$. Within a
period of $512$ there are four representatives of this crossing class, of which one realizes
the full valuation word.
\end{example}

\section{Words with at most 75 odd runs}\label{sec:75}
\begin{lemma}[run-length realizer bound]\label{lem:runbound}
If an odd run starts at $x$ and lasts $\ell$ steps, the endpoint plus one equals
$(3/2)^\ell(x+1)$, so $2^\ell\mid x+1$. At each odd step $x+1$ is multiplied by $3/2$,
and at each even step it decreases. Writing $D_i$ for the number of odd steps before run $i$
and $L=\log_2(n+1)$, we get $2^{\ell_i}\le(n+1)(3/2)^{D_i}$, hence $D_{i+1}\le L+\alpha D_i$ with
$\alpha=\log_23$, and therefore
\[
d\le G_sL,\qquad G_s=\sum_{j=0}^{s-1}\alpha^{\,j}.
\]
\end{lemma}

Neither ingredient is new. The divisibility $2^{\ell}\mid x+1$ across a run of $\ell$ odd
steps appears explicitly as Lemma 8 of Hercher \cite{hercher2023}, written there as
$1\le k_i\le x_i$ with $x_i=\log_2(n_i+1)$; that is the antecedent we are aware of. Garner's
contribution is separate and is recorded below: the coefficient stopping time itself
\cite[\S4]{garner1981} and the Diophantine study of powers of $2$ against powers of $3$
\cite[\S5--\S6]{garner1981}. The geometric sum $G_s=(\alpha^s-1)/(\alpha-1)$ is likewise the sum
$(\delta^{m_2}-1)/(\delta-1)$ appearing in Hercher's Theorem 21, with the \emph{same} constant
$\delta=\log_23=\alpha$; his recurrence $x_{i+1}<\delta x_i$ is the purely geometric form that
cyclicity permits, while the affine form $D_{i+1}\le L+\alpha D_i$ above re-injects the seed
size $L$ at every run because a finite path has a fixed starting point.

The method of the next theorem --- a Diophantine bound on powers of $2$ against powers of $3$,
combined with a finite verification range --- is Garner's \cite[\S6--\S7]{garner1981}; see
Remark \ref{rem:history}.

\begin{theorem}\label{thm:75}
Every eligible realizer of a first-crossing word with at most $75$ odd runs descends at that
crossing. Any counterexample to Terras's conjecture has at least $76$ odd runs, crossing time
$k\ge10\,439\,860\,591$ and odd-input count $d\ge6\,586\,818\,670$.
\end{theorem}

\begin{proof}
Rozier and Terracol \cite[Cor.~5.4]{rozier2026} verify $\chi(n)=\tau(n)$ for $2\le n\le B$ with
$B=28\cdot10^{18}$, so a failure has $n>B$. Put $\Lambda=k\ln2-d\ln3>0$. By
Proposition \ref{prop:comparator} a failure gives $(e^\Lambda-1)n\le M_w$, hence with
\eqref{eq:Mbounds}
\begin{equation}\label{eq:lambda}
0<\Lambda<s/n<s/B,\qquad 0<\Lambda/d<1/(3n)<1/(3B).
\end{equation}

\emph{The two large integers.} Take the convergents of $\alpha=\log_2 3$ of indices
$20$ and $19$, respectively:
\[
 (a,b)=(9\,809\,721\,694,\,6\,189\,245\,291),\qquad
 (u,v)=(630\,138\,897,\,397\,573\,379).
\]
One certifies $a/b<\alpha$ and $(u/v)\ln2-\ln3>1/(3B)$. With the second half of
\eqref{eq:lambda} this forces $a/b<\alpha<k/d<u/v$. The determinant $bu-av=1$ makes $bk-ad$ and
$ud-vk$ positive integers, and expanding gives $(k,d)=(ud-vk)(a,b)+(bk-ad)(u,v)$. Hence
$k\ge a+u=10\,439\,860\,591$ and $d\ge b+v=6\,586\,818\,670$; this pair is the convergent of index
$21$, the mediant of the two used, and it is the pair that occurs in Eliahou's work on nontrivial
cycle lengths \cite{eliahou1993}.

\emph{Excluding at most 75 runs.} Suppose now that $s\le75$, so that $\Lambda<75/B$.
For $p/q<\alpha$ put $e=q\ln3-p\ln2>0$. Since $k/d>\alpha$,
$q\Lambda+de=(qk-pd)\ln2\ge\ln2$. Two pairs suffice, with certified enclosures
$e_{\mathrm{I}}\in(5.22,5.24)\cdot10^{-18}$ at
$(p,q)=(9\,881\,527\,843\,552\,324,\,6\,234\,549\,927\,241\,963)$ and
$e_{\mathrm{II}}\in(1.24,1.25)\cdot10^{-19}$ at
$(p,q)=(6\,724\,555\,128\,221\,608\,268,\,4\,242\,721\,909\,926\,539\,673)$. Pair I gives
$d>(0.69-q\cdot75/B)/(5.24\cdot10^{-18})>1.26\cdot10^{17}$.
If $d<2^{61}<2.31\cdot10^{18}$, pair II gives
$\Lambda>(0.69-2.31\cdot10^{18}\cdot1.25\cdot10^{-19})/(4.25\cdot10^{18})>9\cdot10^{-20}$, hence
$n<75/(9\cdot10^{-20})$ and $n+1<2^{70}$; Lemma \ref{lem:runbound} with $G_{75}<1.72\cdot10^{15}$
then gives $d<70\cdot1.72\cdot10^{15}=1.204\cdot10^{17}$, contradicting $d>1.26\cdot10^{17}$.

\emph{The tail.} For $d\ge2^{61}$, Rhin's effective estimate \cite{rhin1987} (in the form used by
\cite[Prop.~6.3]{rozier2026}) gives $\Lambda\ge k^{-13.3}$, so $n<75(2d)^{13.3}$ and
$\log_2(n+1)<7+14\log_2(2d)$. With $2^j\le d<2^{j+1}$, Lemma \ref{lem:runbound} forces
$2^j\le d<1.72\cdot10^{15}\,[\,7+14(j+2)\,]$.
At $j=61$ the upper bound is $1.52908\cdot10^{18}<2^{61}$. It stays below $2^j$ by induction,
since the bracket increases by $14$ while $2^j$ doubles. This contradiction excludes the remaining
values of $d$.
\end{proof}

\begin{remark}[Dependencies of Theorem \ref{thm:75}]
The Rozier--Terracol verification supplies the seed bound $n>B$. Combined with
\eqref{eq:lambda}, the Farey interval then gives the two integer bounds. Rhin's estimate enters
only in the case $d\ge2^{61}$; its hypothesis is $H=\max(|u_1|,|u_2|)\ge2$
\cite[Prop.~6.3]{rozier2026}, and with $u_1=k$, $u_2=-d$ and $k>d$ one has $H=k$, so the form
used is $\Lambda\ge k^{-13.3}$.
\end{remark}

\begin{remark}[Historical context, and what is claimed here]\label{rem:history}
Every method used above is older than this note; what is new is the combination and the
resulting explicit constants for \emph{finite first crossings}.

\emph{Garner \cite{garner1981}} is the direct precedent. He defined the coefficient stopping
time, also formulated the conjectured equality with the stopping time --- he cites Terras
\cite{terras1976} as his reference [3] --- and verified it for
$n<1\,150\,000$, identified the realizing residue class modulo $2^{1+p(m)}$ with
$p(m)=\lfloor m\log_23\rfloor$, and proved (his \S6) that $\sigma(n)=\kappa(n)$ whenever the
number of odd \emph{terms} $m$ is below a bound assembled from a study of the powers of $2$ and
$3$ --- his tabulated breakpoints are denominators of convergents of $\log_23$ --- together with
a verified range. His \S7 applies this to cycles, obtaining at least $35\,400$ terms from
$n>6\cdot10^7$. Theorem \ref{thm:75} is that argument with two changes: the hypothesis bounds
the number of odd \emph{runs} rather than odd terms, which is strictly weaker because a single
run may contain arbitrarily many odd steps, and the modern verification range and an explicit
Farey sandwich replace his tabulated constants.

\emph{Simons and de Weger \cite{simonsdeweger2005}} developed the run-growth approach for
nontrivial cycles, bounding the number of local minima; \emph{Hercher \cite{hercher2023}}
strengthened that line to $m\ge92$. Their counts and ours are the same kind of quantity but not
the same quantity: in a cycle the odd runs close up, so the number of local minima equals the
number of runs, whereas in a finite first-crossing word the first run may begin at the seed,
which is not a local minimum, and the last run is cut off by the crossing. Their bounds also use
cyclicity, through the closing of the product around the cycle, which a finite path does not
provide. The numerical agreement between the threshold $76$ here and the Simons--de Weger value
$m\ge76$ is a coincidence of similar inputs: with a modern verification range their cycle bound
becomes $m\ge83$ \cite{hercher2023}, while ours stays at $76$. \textbf{Nothing in this note
raises the first-crossing threshold to $92$, and Hercher's cycle conclusion must not be read as
applying to finite first crossings.}

\emph{Lagarias \cite{lagarias1985}} is the closest predecessor of all, and
Remark~\ref{rem:classical} gives the locators. His (2.14) is our descent criterion verbatim
after rearrangement, and his Theorem~E already proves, by exactly our route --- linear forms in
logarithms fed into that criterion --- that for admissible words of length $k\ge k_0$ only the
smallest element of the class can fail. What is left to us is to make $k_0$ explicit under a
hypothesis on the run count, which is what Lemma~\ref{lem:runbound} and
Theorem~\ref{thm:75} do. Lagarias also records \cite[\S2.3]{lagarias1985} that Terras
\cite{terras1976} and ``later Garner'' \cite{garner1981} conjectured the equality, which
settles the question of independence raised above.

\emph{Rozier and Terracol \cite{rozier2026}} supply the verification range and, in their
Remark~5.5, credit Simons with proposing classification by local maxima. Three classes of object
must be kept apart. A \emph{nontrivial cycle} is the subject of
\cite{garner1981,simonsdeweger2005,hercher2023}. A \emph{paradoxical sequence}
\cite[Def.~1.1]{rozier2026} is a finite trajectory with $T^j(n)\ge n$ and coefficient below one
at step $j$, for \emph{any} such $j$. A \emph{first-crossing failure}, the object here, is a
paradoxical sequence with $j=\tau(n)$, hence a proper subclass. Their Theorem 5.3 excludes
paradoxical sequences with $93\le j\le301\,993$ outright, and all with
$4615\le n\le2.8\cdot10^{19}$; the $301\,993$ is the extent of a finite computation, whereas
the bound $k\ge10\,439\,860\,591$ of Theorem \ref{thm:75} is a lower bound on a counterexample
derived from the Farey sandwich. A counterexample must satisfy both, and ours is the stronger
constraint; there is no conflict.

\emph{Hikawa \cite{hikawa2026}} reviews the problem and singles out exactly the triple studied
here --- the correction term $q_v$, the coefficient difference $2^k-3^d$, and the least positive
representative $r_v$ of the residue class. \textbf{Priority for the specific first-crossing
statements of this note is therefore not established}; we have not located a match, which is not
the same as there being none.
\end{remark}

\section{A family with arbitrarily many runs}
\begin{proposition}\label{prop:110}
Repeat the block $110$ exactly $h\ge1$ times and append zeros up to the first coefficient
crossing. Every eligible realizer descends at that crossing.
\end{proposition}
\begin{proof}
One block acts by $H(x)+5=9(x+5)/8$, so $-5$ is its fixed point.
The block's full coefficient is $9/8$ and its two nonempty proper-prefix coefficients
are $3/2$ and $9/4$. Consequently the first crossing occurs among the appended zeros
at $k=\lceil h\log_2 9\rceil\le4h$. Iteration gives
$H^h(n)+5=9^h(n+5)/8^h$. Integrality implies $8^h\mid n+5$, hence $n\ge8^h-5$.
The normalized memory is $5(1-(8/9)^h)<5$ and the appended zeros leave it unchanged.

For $h\le63$, $k\le252$ is below the crossing-length bound in Theorem \ref{thm:75}.
For $h\ge64$, Rhin's estimate gives $n\Lambda\ge(8^h-5)/(4h)^{13.3}>5$.
Indeed, $8^{64}-5>5\cdot256^{14}$, and $((h+1)/h)^{14}\le(65/64)^{14}<8$,
so $8^h-5>5(4h)^{14}$ persists by induction. A failure requires
$n\Lambda<n(e^\Lambda-1)\le M_w<5$, a contradiction.
\end{proof}
The prefix $(110)^h$ and the completed word generally have different least seeds.
Write $m=k-3h$ and $n=8^h t-5$. All these $n$ form the binary prefix class,
since $2^{3h}=8^h$ is its exact modulus. Its least positive representative is $8^h-5$.
To append $m$ zeros we also need $2^m\mid 9^h t-5$, the block endpoint.
At $h=12$ this gives $k=39$, $m=3$ and $t\equiv5\pmod8$. The prefix and completed-word
least seeds are therefore $68\,719\,476\,731$ and $343\,597\,383\,675$, respectively.

\section{A finite-word repetition bound}
\begin{proposition}\label{prop:square}
Let a first-crossing word from $n$ begin with $UU$, where $|U|=\ell\ge1$.
Then $(4/3)^\ell<n+1$. If the crossing fails to give descent, then
\[
 \ell<\frac{\log_2(n+1)}{\log_2(4/3)}
      <\frac{14.3}{\log_2(4/3)}\log_2 k<35\log_2 k .
\]
\end{proposition}
\begin{proof}
Set $x=T^\ell(n)$. Both $n$ and $x$ start with $U$, so the parity-class congruence gives
$x\equiv n\pmod{2^\ell}$. Since $\ell<k$, $\rho_\ell\ge1$.
Equality would give $3^{d_\ell}=2^\ell$, impossible for $\ell\ge1$.
Thus $x\ge\rho_\ell n>n$, so $x-n$ is a positive multiple of $2^\ell$.
Equivalently $2^\ell(x-n)$ is a positive multiple of $2^{2\ell}$.
At each step $T(y)+1\le(3/2)(y+1)$, whence
\[
 2^\ell\le x-n<x+1\le(3/2)^\ell(n+1).
\]
Dividing gives the first assertion, independently of whether the seed descends.

For a failure set $\Lambda=k\ln2-d\ln3>0$. Then
$n\Lambda<n(e^\Lambda-1)\le M_w\le d/3$ by Proposition~\ref{prop:comparator}
and \eqref{eq:Mbounds}. Since $k\ge2$ and $d\le k$, Rhin's bound
\cite[Prop.~6.3]{rozier2026} gives $\Lambda\ge k^{-13.3}$. It follows that
$n<(d/3)k^{13.3}\le k^{14.3}/3$ and hence $n+1<k^{14.3}$.
Taking logarithms completes the proof, since $14.3/\log_2(4/3)<35$.
\end{proof}
The restriction concerns repetition inside a single finite first-crossing prefix.
The periodic-approximation results in \cite{dj-sturmian2026} concern infinite parity words.

\section{A restriction from reciprocal sums}
Write $x_j=T^j(n)$ and form the product over the odd input values of the first $k$ steps,
\[
 Q=\prod_{\substack{0\le j<k\\x_j\ {\rm odd}}}\left(1+\frac{1}{3x_j}\right).
\]
Multiplying the single-step factors gives $x_k/n=\rho_k Q$. Hence
Proposition~\ref{prop:comparator} gives $Q=1+M_w/n$.

Let $D=28\,019\,077\,177\,231\,758\,495$. For any finite shortcut path with distinct
inputs all exceeding $D$, let $R$ be the sum of the reciprocals of its odd inputs.
Nikzad, Frinkle and Bagherzadeh prove $R<B_0$ in \cite[Theorem~1A.1]{nfb2026}, where
\[
 B_0=\frac{45}{256}+10^{-6}+10^{-15}.
\]
Their proof gives a second rational upper bound by retaining the coefficient in (153)
before the rounding in (154). Set
\[
 K=\frac{6\,719\,267}{13\,693\,680},\qquad
 c=\frac{93\,657\,291\,203\,213}{10^{15}},\qquad
 B_1=\frac{15}{4}\,Kc+10^{-6}+10^{-15}.
\]
Indeed, for each short gap bracket of total length $m$, (153) bounds the sum over
its admissible starting values by $K2^{-m}$. The one-step balance bound (142) contributes
the factor $15/4$. Equation (164) bounds the weighted bracket count by the exact dyadic
$\sum_{m=61}^{1000}\widehat C_m2^{-m}=0.09365729120321290556\ldots<c$.
Adding the long-shell allowance from (139) and the boundary allowance from (149),
as in (151)--(154), gives $R<B_1<B_0$. These equation numbers refer to \cite{nfb2026}.
The unpublished verification \cite{dj-indep2026} independently reconstructs the high-path
estimate and reports an exact reproduction of the finite count in (164).

\begin{proposition}\label{prop:recip}
If a counterexample exists, let $n$ have the least crossing length $k$ among all
counterexamples. For $i=0,1$ it satisfies
\[
 0<\Lambda=-\ln\rho_k < B_i/3,\qquad
 \rho_k>\exp(-B_i/3).
\]
\end{proposition}
\begin{proof}
We use the external height record $H(y)\le2456$ through $D$, reported in
\cite[eq.~(34)]{nfb2026}. Here $H$ counts ordinary Collatz operations to first arrival
at $1$. For $2\le y\le D$, this gives $\tau(y)\le\chi(y)\le H(y)\le2456$, since
shortcut steps combine ordinary operations. The lower bound on $k$ in Theorem~\ref{thm:75} therefore gives $n>D$.
For $0\le j<k$, the prefix coefficient is at least one and its affine correction
is nonnegative, so $x_j\ge n>D$.

If two of $x_0,\ldots,x_{k-1}$ coincided, the intervening segment would contain a
cycle of length $L<k$. Let $m$ be its minimum. One full turn from $m$ has positive
affine correction, since any positive cycle has an odd input. The return to $m$
therefore forces its coefficient below one. Thus $\tau(m)\le L$, while every
value on the cycle is at least $m$, giving a failure of shorter crossing length.
This proves that the $k$ inputs are distinct. The high-path estimate applies even
when the endpoint repeats an input. Finally $x_k\ge n$ gives $1\le\rho_k Q$, and
\[
 \Lambda\le\ln Q<\frac13\sum_{\substack{0\le j<k\\x_j\ {\rm odd}}}\frac1{x_j}
 < B_i/3.
\]
\end{proof}

In the coordinate $-\log_2\rho_k$, the nominal first-crossing interval has length
one. For each bound the surviving interval has length $B_i/(3\ln2)$. The values
below are rounded for display.
\[
\begin{array}{c|c|c|c}
 &\exp(-B_i/3)&\text{surviving width}&\text{excluded width}\\\hline
 B_0&0.9430895072&8.45334\%&91.54666\%\\
 B_1&0.9441733669&8.28763\%&91.71237\%
\end{array}
\]
These percentages measure intervals of log coefficients. The restriction concerns a
counterexample of minimal crossing length.

\section{The open question}\label{sec:open}
A shortest counterexample would have at least $76$ odd runs, obey the coefficient bound in
Proposition \ref{prop:recip}, and have no initial square whose period exceeds the bound in
Proposition \ref{prop:square}. For each remaining first-crossing word $w$, its coefficient
and its additive correction are known exactly. The open estimate for its least eligible
realizer $r(w)$ is
\[
             r(w)>\frac{C_w}{2^k-3^d}=\frac{M_w}{\rho_k^{-1}-1}.
\]
Equality already gives a counterexample. A lower bound for $r(w)$ or an upper bound
for the right side can certify a word if it preserves this strict inequality.
The rational-shadow estimates of \cite{dj-shadow2026} supply bounds for individual blocks.
Applying one to this question requires transporting it back to the starting seed and comparing
it with the threshold of the actual first-crossing prefix.

\newpage
\section*{Status of the results}
It is worth separating three things explicitly.

\emph{Established before this note.} The residue-class correspondence for parity words
(Terras \cite{terras1976}); the affine form, the descent criterion
$n>\varepsilon_k^{-1}\rho_k(v)$, and --- in Theorem~E --- the statement that for admissible
words of length $k\ge k_0$ only the smallest element of a class can fail, proved from linear
forms in logarithms (Lagarias \cite[\S2.2--\S2.3, Thm.~E]{lagarias1985}); the coefficient stopping time, the conjecture
$\kappa=\sigma$, and the method of combining a Diophantine bound on $2^a$ against $3^b$ with a
finite verification range, including a descent theorem under a bound on the number of odd
\emph{terms} (Garner \cite{garner1981}); run-growth bounds for the number of local minima in a
nontrivial cycle (Simons--de Weger \cite{simonsdeweger2005}, Hercher \cite{hercher2023});
verification of $t=\tau$ to $2.8\cdot10^{19}$ and the exclusion of paradoxical sequences with
$93\le j\le301\,993$ (Rozier--Terracol \cite{rozier2026}); the reciprocal-sum estimate
(Nikzad--Frinkle--Bagherzadeh \cite{nfb2026}).

\emph{Offered here.} The transfer of Garner's method to finite first crossings with the
hypothesis weakened from odd \emph{terms} to odd \emph{runs}, which is what
Lemma \ref{lem:runbound} buys; the resulting explicit constants
$k\ge10\,439\,860\,591$, $d\ge6\,586\,818\,670$ and a threshold of $76$ runs; descent for
the whole repeated-$110$ family, which has arbitrarily many runs; the finite-word square bound
$|U|<35\log_2k$; the effective and non-asymptotic form of Lagarias's Theorem~E that these constants provide;
and the
reciprocal-sum restriction $\rho_k>0.9441733669$ on a shortest counterexample. Subject to the
priority caveat of Remark \ref{rem:history}.

\emph{Unresolved.} A lower bound on the least positive realizer of a many-run word with a
shallow crossing and no long initial repetition --- the inequality of
\S\ref{sec:open} --- is open. Whether Hercher's cycle techniques transfer to finite
first-crossing trajectories, and in particular whether the threshold of $76$ runs can be raised,
is left for future work; \textbf{no such raising is attempted or claimed here}. Terras's
conjecture itself is open, as is the Collatz conjecture.

\section*{Reproducibility}
The companion verifier checks the numerical inequalities used in the proofs,
recomputes the finite bracket count, and verifies the worked examples using integer and
rational arithmetic. It lives in the
companion repository
\href{https://github.com/innerlightr-wq/terras-first-crossing}{\texttt{github.com/innerlightr-wq/\allowbreak terras-first-crossing}}, with the verification
code pinned at the immutable commit
\href{https://github.com/innerlightr-wq/terras-first-crossing/commit/a93b1e82d13fab7024c380ae36320e4ec62a73f9}{\texttt{a93b1e8}}.
Run the following command from the project root.
\begin{verbatim}
python3 scripts/verify_crossing.py
\end{verbatim}
The script checks the numerical inequalities with integers and rationals,
encloses $\ln2$ and $\ln3$ by a 96-term series with a geometric tail, recomputes the
finite bracket sum in (164) of \cite{nfb2026}, and encloses each exponential cutoff by
alternating Taylor bounds. The SHA-256 digest of the verifier begins
\texttt{f852c88311071a0d}. The orbit and residue-class examples are finite checks.
The proofs for arbitrary word lengths are in the text.

The finite seed verification of Rozier--Terracol, Rhin's estimate, the high-path
accounting of Nikzad--Frinkle--Bagherzadeh and the height record cited above are
external inputs. The general case of Terras's conjecture remains open.

\section*{AI assistance and author responsibility}
The author directed this research and used Claude (Anthropic) and ChatGPT (OpenAI) for
proof exploration and revision, cross-checking the reciprocal-sum estimates against the
source paper, developing and testing the exact-arithmetic verifier, and drafting the
manuscript. The author takes responsibility for the mathematical claims, their
interpretation, and any errors.

\begin{thebibliography}{10}

\bibitem{eliahou1993}
Shalom Eliahou.
\newblock The $3x+1$ problem: new lower bounds on nontrivial cycle lengths.
\newblock {\em Discrete Mathematics}, 118:45--56, 1993.

\bibitem{garner1981}
Lynn~E. Garner.
\newblock On the {C}ollatz $3n+1$ algorithm.
\newblock {\em Proceedings of the American Mathematical Society}, 82(1):19--22,
  1981.

\bibitem{hercher2023}
Christian Hercher.
\newblock There are no {C}ollatz-$m$-cycles with $m\le91$.
\newblock {\em Journal of Integer Sequences}, 26(3), 2023.
\newblock Article 23.3.5. Corrigendum (Theorem 21),
  \url{https://cs.uwaterloo.ca/journals/JIS/VOL26/Hercher/corrigendum.pdf}; the
  conclusion $m\ge92$ is unaffected.

\bibitem{hikawa2026}
Kazunobu Hikawa.
\newblock Research notes on {Terras}' coefficient stopping time problem {I}:
  Previous results and open problems, 2026.
\newblock English translation, November 2026; Japanese original, Zenodo, 13
  November 2026, DOI 10.5281/zenodo.24187839.

\bibitem{dj-sturmian2026}
Elias~De Jes\'us.
\newblock Beyond {S}turmian words in the $3x+1$ periodicity conjecture: A
  periodic-approximation and arithmetic-height bridge, 2026.
\newblock Revision 5, 13 November 2026; Revision 4 deposited as Zenodo DOI
  10.5281/zenodo.23070093.

\bibitem{dj-shadow2026}
Elias~De Jes\'us.
\newblock Exact realizer congruences and rational-shadow bounds for {C}ollatz
  words, 2026.
\newblock Zenodo preprint, 31 October 2026.
  \href{https://doi.org/10.5281/zenodo.23847289}{\texttt{doi:10.5281/zenodo.23847289}}.

\bibitem{dj-floors2026}
Elias~De Jes\'us.
\newblock Exact realizer floors for one- and two-defect valuation words in the
  accelerated $3x+1$ map, 2026.
\newblock Zenodo preprint, 22 October 2026.
  \href{https://doi.org/10.5281/zenodo.23579193}{\texttt{doi:10.5281/zenodo.23579193}}.

\bibitem{dj-indep2026}
Elias~De Jes\'us.
\newblock Independent verification of high-path reciprocal bounds for
  nonrepeating {C}ollatz trajectories, 2026.
\newblock Unpublished manuscript, 3 November 2026.

\bibitem{lagarias1985}
Jeffrey~C. Lagarias.
\newblock The $3x+1$ problem and its generalizations.
\newblock {\em American Mathematical Monthly}, 92(1):3--23, 1985.

\bibitem{nfb2026}
Behzad Nikzad, Karl Frinkle, and Fatemeh Bagherzadeh.
\newblock On bounds for reciprocal sums and residues in {C}ollatz orbits, and
  their consequences, 2026.
\newblock Hexagon preprint 2611.00033v1,
  \url{https://hexagonmath.org/2611.00033v1}.

\bibitem{rhin1987}
Georges Rhin.
\newblock Approximants de {P}ad\'e et mesures effectives d'irrationalit\'e.
\newblock In {\em S\'eminaire de Th\'eorie des Nombres, Paris 1985--86},
  volume~71 of {\em Progress in Mathematics}, pages 155--164. Birkh\"auser
  Boston, 1987.

\bibitem{rozier2026}
Olivier Rozier and Claude Terracol.
\newblock Paradoxical behavior in {C}ollatz sequences.
\newblock {\em Discrete Mathematics}, 349:115167, 2026.
\newblock arXiv:2502.00948v5.

\bibitem{simonsdeweger2005}
John Simons and Benne de~Weger.
\newblock Theoretical and computational bounds for $m$-cycles of the $3n+1$
  problem.
\newblock {\em Acta Arithmetica}, 117(1):51--70, 2005.
\newblock The bound was subsequently raised; Hercher \cite{hercher2023} reports
  $m\ge76$ for Simons--de Weger, and $m\ge83$ on substituting a modern
  verification range.

\bibitem{terras1976}
Riho Terras.
\newblock A stopping time problem on the positive integers.
\newblock {\em Acta Arithmetica}, 30:241--252, 1976.

\end{thebibliography}
\end{document}
